\begin{sectionbox}
	\subsection{Stefan Boltzmann}
	\begin{emphbox}
		$P=\epsilon \sigma A T^4 = 4 \pi r^2 I \qquad [P]=W$
	\end{emphbox}
	mit $A$ bestrahlte Fläche und $T$ Temperatur. \\
	Für schwarze Körper gilt: $\epsilon = 1$\\
	Wienscher Verschiebungssatz: $\lambda_{\ir max}=\frac{\SI{2,898e-3}{\kelvin}}{T} \si{\meter}$ \\
	Größtmögliche Wellenlänge bei einer Temperatur \\

	\textbf{Lichtintensität} $I=\frac{P}{4\pi r^2}$
\end{sectionbox}